\subsection{多线性态射}

7.3.1. 设 \(M_{1},\ldots,M_{d}\) 和 N 是以 B 为基底的向量丛，u 是从集合 \(M_{1}\times_{B}\ldots\times_{B}M_{d}\) 到 N 的映射。如果满足以下条件，则称 u 为多线性态射（或 d-线性态射）：

对于每个 \(b_0 \in \mathbf{B}\)，存在一个开邻域 \(\mathbf{U}\)，包含 \(b_0\) 且位于 \(\mathbf{B}\) 中，存在向量丛图 \(t^j = (\mathbf{U}, \varphi^j, \mathbf{F}^j)\)，它们属于 \(\mathbf{M}_j\)（\(1 \leqslant j \leqslant d\)），以及向量丛图 \(t = (\mathbf{U}, \varphi, \mathbf{F})\)，它属于 \(\mathbf{N}\)，还有映射 \(\lambda\)，其类为 \(\mathbf{C}^r\)，从 \(\mathbf{U}\) 到巴拿赫空间 \(\mathcal{L}(\mathbf{F}^1, \ldots, \mathbf{F}^d; \mathbf{F})\)，该空间由从 \(\mathbf{F}^1 \times \cdots \times \mathbf{F}^d\) 到 \(\mathbf{F}\) 的连续 d-线性映射组成，使得：

\[
(t _ {b} \circ \lambda (b)) (x _ {1}, \dots , x _ {d}) = u (t _ {b} ^ {1} (x _ {1}), \dots , t _ {b} ^ {d} (x _ {d}))
\]

对所有 \(b \in \mathbf{U}\) 和所有 \(x_{j} \in \mathbf{F}^{j}\) 都成立。

每个多线性态射 \(u\) 都是从纤维积 \(\mathbf{M}_{1} \times_{\mathbf{B}} \cdots \times_{\mathbf{B}} \mathbf{M}_{d}\) 到 \(\mathbf{N}\) 的流形态射，并且对每个 \(b \in \mathbf{B}\) 诱导出一个连续 \(d\)-线性映射 \(u_b\)，从 \((\mathbf{M}_{1})_{b} \times \cdots \times (\mathbf{M}_{d})_{b}\) 到 \(\mathbf{N}_{b}\)。

若 \(f\) 是从 B 到流形 \(\mathbf{B}'\) 的态射，则从 \(\mathbf{M}_1 \times_{\mathbf{B}} \cdots \times_{\mathbf{B}} \mathbf{M}_d\) 到向量丛 \(\mathbf{M}'\)（其基底为 \(\mathbf{B}'\)）的多线性态射，按定义是从 \(\mathbf{M}_1 \times_{\mathbf{B}} \cdots \times_{\mathbf{B}} \mathbf{M}_d\) 到 \(f^*\mathbf{M}'\) 的多线性态射与由 \(f\) 给出的从 \(f^*\mathbf{M}'\) 到 \(\mathbf{M}'\) 的典范态射的复合。

双线性态射也称为配对。当 \(d=1\) 时，线性态射就是 7.2.1 意义下的态射。当 \(d=0\) 时，0-线性态射与 N 的一个截面（7.4）等同。

7.3.2. 以 B 为基底的代数丛是以 B 为基底、配备了从 \(A \times_{B} A\) 到 A 的配对的向量丛 A。于是每个纤维 \(A_{b}\) 都带有一个 K-代数结构。若对每个 \(b \in B\)，代数 \(A_{b}\) 有单位元，记为 \(e_{b}\)，则映射 \(b \mapsto e_{b}\) 是 A 的一个截面（参见 7.4）。如果对 B 中每个点 \(b_{0}\)，都存在一个向量丛图 \(t = (\mathrm{U}, \varphi, \mathrm{E})\)，它属于 A 并在点 \(b_{0}\) 处，以及 E 上的 K-代数结构，使得 \(t_{b}\) 是从 E 到 \(A_{b}\) 的代数同构，对所有 \(b \in U\) 均如此，则称代数丛 A 是局部平凡的。

7.3.3. 设 A 是以 B 为基底的带单位元的结合代数丛。以 B 为基底的 A-模丛定义为以 B 为基底的向量丛 M，配备一个配对 \(m: A \times_{B} M \to M\)，使得映射 \(m_{b}: A_{b} \times M_{b} \to M_{b}\) 对每个 \(b \in B\) 定义出 \(A_{b}\) 在纤维 \(M_{b}\) 上的模结构。

设 M 和 M′ 是两个 A-模向量丛。从 M 到 M′ 的 A-同态定义为以 B 为基底的向量丛态射 \(g: M \rightarrow M'\)，它对 B 中每个 b 都诱导出一个 \(A_b\)-线性映射，从 \(M_b\) 到 \(M_b'\)。

设 A 是局部平凡的。若对 B 中每个点 \(b_0\)，都存在一个向量丛图 \(t = (U, \varphi, E)\)，它是如 7.3.2 中所述的 A 在 \(b_0\) 处的图；另有一个向量丛图 \(t' = (U, \varphi', L)\)，它是 M 在 \(b_0\) 处的图；以及 L 上的 E-模结构，使得对每个 \(b \in U\)，\(t_b'\) 是 \(t_b\)-同构，将 \(A_b\)-模 \(M_b\) 映到 E-模 L，则称 A-模丛 M 是局部平凡的。

7.3.4. 设 A 是 K 上的巴拿赫代数（例如，带有有限维 K-代数结构的域）。平凡丛 \(A_{B}\) 于是是一个局部平凡代数丛。A-模 \(A_{B}\) 中的局部平凡丛 M 也称为 A 上的向量丛（以 B 为基底）。于是纤维 \(M_{b}\) 是拓扑 A-模。从 M 到另一个 A 上的向量丛的 f-态射 u 称为 A-线性的，如果映射 \(u_{b}\) 对所有 \(b \in B\) 都是 A-线性的。

